\documentclass[12pt]{article}
\usepackage{amsmath}
\usepackage{amssymb}
\usepackage[margin=1in]{geometry}
\usepackage{hyperref}

\title{Methods: Analysis and Visualization Pipeline}
\date{}

\begin{document}
\maketitle

\section{Disk Identification}

All non-trivial analyses depend on a two-stage disk identification algorithm
applied to each snapshot independently.

\subsection*{Center and reference frame}

The disk center is taken to be the position of the most massive sink particle
(PartType5). Before any sinks form, the center is placed at the densest gas
particle within a fiducial search radius of the Jeans-refinement region. Using
the most massive sink (rather than the mass-weighted center of all sinks) avoids
frame jumps when secondary fragments form far from the primary disk. The bulk
velocity of the system is subtracted by computing the mass-weighted mean velocity
of all gas within a search radius $r_\mathrm{search} = 10\,\mathrm{AU}$ of the
center.

\subsection*{Disk orientation}

The disk angular momentum axis $\hat{L}$ is computed from all gas particles
within $r_\mathrm{search}$:
\begin{equation}
    \hat{L} = \frac{\sum_i m_i\,(\mathbf{r}_i \times \mathbf{v}_i)}
                   {\left|\sum_i m_i\,(\mathbf{r}_i \times \mathbf{v}_i)\right|}
\end{equation}
where $\mathbf{r}_i$ and $\mathbf{v}_i$ are positions and velocities in the
center-of-mass frame. A rotation matrix is constructed to align $\hat{L}$ with
$\hat{z}$, transforming all particle coordinates and velocities into the disk
frame.

\subsection*{Keplerian velocity}

The enclosed mass for each gas particle is computed by sorting particles by
cylindrical radius $r_\mathrm{cyl}$ and taking the exclusive cumulative sum of
gas masses plus the stellar contribution interior to each particle's orbit. The
Keplerian speed is then
\begin{equation}
    v_K = \sqrt{\frac{G\,M_\mathrm{enc}(<r_\mathrm{cyl})}{r_\mathrm{cyl}}}
\end{equation}

\subsection*{Selection criteria}

A particle is classified as disk gas if it simultaneously satisfies:
\begin{enumerate}
    \item $r_\mathrm{cyl} < r_\mathrm{max} = 2\,\mathrm{kAU}$ \quad (spatial)
    \item $|z|/r_\mathrm{cyl} < 0.3$ \quad (geometric flatness)
    \item $\rho > 10^{-15}\,\mathrm{g\,cm^{-3}}$ \quad (density threshold)
    \item $v_\phi > 0$ \quad (co-rotating with the disk)
    \item $v_\phi/v_K > 0.3$ \quad (at least 30\% Keplerian)
\end{enumerate}
all evaluated in the rotated disk frame. The resulting kinematically selected set
is then extended to its bounding cylinder (100th percentile in $r_\mathrm{cyl}$
and $|z|$) to recover any disk material just outside the strict kinematic
boundary.

% ─────────────────────────────────────────────────────────────────────────────

\section{Surface Density Maps}

Face-on and edge-on surface density maps are produced using the \texttt{Meshoid}
library, which evaluates the SPH kernel-weighted projection of the gas mass onto
a uniform pixel grid. The surface density at pixel $(x,y)$ is
\begin{equation}
    \Sigma(x,y) = \sum_i m_i\,W\!\left(|\mathbf{r}_{i,\perp} - \mathbf{r}_\mathrm{pix}|,\,h_i\right)
\end{equation}
where $W$ is the cubic spline kernel and $h_i$ is the SPH smoothing length of
particle $i$. Maps are rendered at $400\times400$ pixels over a $4\,\mathrm{kAU}$
full-width box in the face-on frame. A second face-on panel is rendered over a
$40\,\mathrm{kAU}$ box to show the wider environment. Edge-on panels use the same
$4\,\mathrm{kAU}$ width with $y$ and $z$ axes swapped so the disk midplane
appears horizontal. All four panels share a single colorbar spanning
$10^5$--$10^8\,M_\odot\,\mathrm{pc}^{-2}$, fixed globally across all snapshots
to allow direct comparison.

% ─────────────────────────────────────────────────────────────────────────────

\section{Turbulent Velocity Maps and Streaming Subtraction}

To isolate the turbulent component of the velocity field, mean streaming motions
are subtracted from each particle before projection. Gas velocities in the disk
frame are decomposed into cylindrical components $(v_r, v_\phi, v_z)$. The mean
radial and azimuthal streaming profiles $\bar{v}_r(r)$ and $\bar{v}_\phi(r)$ are
computed as mass-weighted averages in 20 annular bins out to the 95th percentile
of the particle distribution. The residual turbulent velocity of each particle is
\begin{equation}
    |\delta v_i| = \sqrt{(\delta v_r)^2 + (\delta v_\phi)^2 + v_z^2}
\end{equation}
where $\delta v_r = v_{r,i} - \bar{v}_r(r_i)$ and
$\delta v_\phi = v_{\phi,i} - \bar{v}_\phi(r_i)$.

Face-on maps of the mass-weighted mean $\langle|\delta v|\rangle$ and its
dispersion $\sigma_{|\delta v|}$ are then computed via Meshoid projection:
\begin{equation}
    \langle|\delta v|\rangle(x,y)
    = \frac{\sum_i m_i\,|\delta v_i|\,W}{\sum_i m_i\,W},
    \qquad
    \sigma_{|\delta v|}^2
    = \frac{\sum_i m_i\,|\delta v_i|^2\,W}{\sum_i m_i\,W}
      - \langle|\delta v|\rangle^2
\end{equation}
The same subtraction is applied to the edge-on projection. Colorbars are fixed
globally across all snapshots by a pre-scan that computes the 99th percentile of
$|\delta v|$ over all frames.

% ─────────────────────────────────────────────────────────────────────────────

\section{Toomre $Q$ Analysis}

\subsection*{Azimuthally averaged profile}

For each snapshot the gas is binned into 20 annular bins in the face-on frame.
In each annulus $b$ we compute mass-weighted quantities:
\begin{equation}
    \Sigma_b = \frac{\sum_{i\in b} m_i}{\pi(r_{\mathrm{hi},b}^2 - r_{\mathrm{lo},b}^2)},
    \qquad
    \sigma_{r,b} = \sqrt{\frac{\sum_{i\in b} m_i\,v_{r,i}^2}{\sum_{i\in b} m_i} - \bar{v}_{r,b}^2}
\end{equation}
The angular velocity is measured directly from the rotation profile:
\begin{equation}
    \Omega_b = \frac{|\bar{v}_{\phi,b}|}{r_b}
\end{equation}
where $\bar{v}_{\phi,b}$ is the mass-weighted mean azimuthal velocity in annulus
$b$ and $r_b$ is the bin center. The epicyclic frequency is computed numerically
from the measured rotation curve via
\begin{equation}
    \kappa_b^2 = \frac{2\,\Omega_b}{r_b}\,\frac{\mathrm{d}(r^2\,\Omega)}{\mathrm{d}r}\bigg|_{r_b}
\end{equation}
where the derivative is evaluated using central differences (one-sided at bin
edges). For a Keplerian rotation curve ($\Omega \propto r^{-3/2}$) this gives
$\kappa = \Omega$, while for a flat rotation curve ($v_\phi = \mathrm{const}$,
$\Omega \propto r^{-1}$) one obtains $\kappa = \sqrt{2}\,\Omega$. The Toomre
parameter per annulus is then
\begin{equation}
    Q_b = \frac{\sigma_{r,b}\,\kappa_b}{\pi\,G\,\Sigma_b}
\end{equation}

\subsection*{Face-on $Q$ map}

A pixel-by-pixel face-on $Q$ map is produced by computing $\sigma_r(x,y)$ as a
Meshoid-projected mass-weighted radial velocity dispersion map, interpolating the
azimuthally averaged $\kappa(r)$ profile onto the pixel grid, and combining with
the Meshoid surface density map $\Sigma(x,y)$:
\begin{equation}
    Q(x,y) = \frac{\sigma_r(x,y)\;\kappa(r_{x,y})}{\pi\,G\,\Sigma(x,y)}
\end{equation}
where $r_{x,y} = \sqrt{x^2 + y^2}$.
The $Q=1$ stability boundary is drawn as a contour.

\subsection*{Space--time heatmap}

Per-snapshot azimuthally averaged $Q(r)$ profiles are saved and assembled into a
$Q(r,t)$ heatmap with $r$ on the $y$-axis and $t - t_1$ on the $x$-axis, where
$t_1$ is the formation time of the first sink particle read directly from the
\texttt{StellarFormationTime} field in the HDF5 output. Sink formation events are
marked as gold stars at their birth radii, with positions determined by
cross-snapshot sink tracking (using the most massive sink as the centering
reference for consistency with position history trajectories). Sinks forming
beyond the plotted $r$-range are omitted, and the primary sink (born at $r = 0$)
is plotted at the lower $r$ boundary.

% ─────────────────────────────────────────────────────────────────────────────

\section{Radial Profiles}

All profiles are mass-weighted averages over annular bins (20 bins for velocity
quantities; 100 bins for the density profile, to better resolve the steep inner
gradient).

\paragraph{Volumetric density $\rho(r)$.}
The GIZMO density field is converted from code units to CGS and mass-weighted per
annulus. Plotted on log--log axes.

\paragraph{Sound speed $c_s(r)$.}
Derived from the stored specific internal energy $u$ via
\begin{equation}
    c_s = \sqrt{\gamma(\gamma - 1)\,u}, \qquad \gamma = 5/3
\end{equation}
where $u$ is in units of $(\mathrm{km\,s^{-1}})^2$.

\paragraph{Radial velocity dispersion $\sigma_r(r)$.}
Computed as in the Toomre analysis (Section~4).

\paragraph{Turbulent speed $\langle|\delta v|\rangle(r)$.}
Mass-weighted mean of the streaming-subtracted $|\delta v_i|$ per annulus (see
Section~3).

\paragraph{Sonic Mach number $\mathcal{M}(r)$.}
\begin{equation}
    \mathcal{M}(r) = \frac{\langle|\delta v|\rangle(r)}{c_s(r)}
\end{equation}
Plotted on a logarithmic $y$-axis with the $\mathcal{M} = 1$ line marking the
sonic transition.

% ─────────────────────────────────────────────────────────────────────────────

\section{Phase Diagrams}

Two-panel phase diagrams are produced for each snapshot. Gas temperature is
derived from the GIZMO specific internal energy:
\begin{equation}
    T = \frac{(\gamma - 1)\,u\,m_p\,\mu}{k_B}
\end{equation}
where $m_p$ is the proton mass and $\mu = 1$ (primordial neutral hydrogen
approximation). Each panel is a 2D mass-weighted histogram in logarithmic bins
($150 \times 150$):
\begin{itemize}
    \item Left panel: $T$ vs.\ $\rho$
    \item Right panel: $f_{\mathrm{H}_2}$ vs.\ $\rho$, where the molecular
          hydrogen fraction $f_{\mathrm{H}_2}$ is read directly from the
          simulation output
\end{itemize}
Disk-classified particles (Section~1) are overlaid as contours to separate the
disk thermodynamic state from the surrounding medium. Colorbars are fixed globally
across all snapshots by a pre-scan to allow direct comparison of the evolving
phase structure.

% ─────────────────────────────────────────────────────────────────────────────

\section{Mass Evolution and Accretion Law}

\paragraph{Disk gas mass.}
$M_\mathrm{disk}(t)$ is the sum of masses of all gas particles identified as disk
members by the algorithm in Section~1. A complementary fixed-aperture estimate
uses a sphere of radius $5\,r_\mathrm{max} = 10\,\mathrm{kAU}$ centered on the
primary sink.

\paragraph{Stellar accretion rate.}
$\dot{M}_*(t)$ is estimated as the finite difference of the total sink mass
between consecutive snapshots divided by the elapsed time.

\paragraph{Accretion law.}
A power-law $M_*(t) = A\,(t - t_1)^\alpha$ is fit to the total sink mass for all
snapshots with $t > t_1$ and $M_* > 1\,M_\odot$, using a linear fit in log--log
space. A reference line $M_* \propto t^2$ (constant accretion rate onto a growing
star) is overlaid for comparison.

\paragraph{Sphere of influence.}
The gravitational sphere of influence $r_\mathrm{SOI}(t)$ is defined as the
3D radius at which the enclosed gas mass equals the total stellar mass. It is
computed by sorting all gas particles by 3D distance from the primary sink,
forming the cumulative gas mass profile $M_\mathrm{gas,enc}(r)$, and solving
$M_\mathrm{gas,enc}(r_\mathrm{SOI}) = M_*(t)$ by linear interpolation.

% ─────────────────────────────────────────────────────────────────────────────

\section{Velocity Power Spectrum}

The three-dimensional velocity power spectrum is computed for each snapshot using
the \texttt{GetPowerSpectrum} routine from the \texttt{starforge\_tools} library
\cite{grudic2021}, adapted to our disk-frame analysis.

\subsection*{Streaming subtraction}

Streaming subtraction is performed as described in Section~3, yielding a
per-particle residual turbulent velocity vector
$\delta\mathbf{v}_i = (\delta v_x,\,\delta v_y,\,v_z)$ in the face-on Cartesian
frame. The vertical component $v_z$ requires no mean subtraction because there is
no mean vertical streaming in a flat disk after COM velocity removal.

\subsection*{Volumetric gridding}

The three components of $\delta\mathbf{v}$ are interpolated onto a $128^3$
uniform grid using Meshoid's 3D SPH kernel interpolation (\texttt{InterpToGrid})
over the disk box of full width $L = 4\,\mathrm{kAU}$, giving a voxel scale of
$\Delta x = L/128 \approx 31\,\mathrm{AU}$.

\subsection*{Power spectrum}

The 3D FFT of each velocity component grid is computed via
\texttt{scipy.fft.fftn}. The total power per mode is
$|V_\mathbf{k}|^2 = |V_{k,x}|^2 + |V_{k,y}|^2 + |V_{k,z}|^2$. Modes are
sorted into spherical shells by their integer wavenumber magnitude
$k_\mathrm{int} = |\mathbf{n}|$ where
$\mathbf{n} \in \mathbb{Z}^3$ is the rounded 3D frequency index. The power
spectral density per shell is
\begin{equation}
    E(k_\mathrm{int}) = \frac{\sum_{|\mathbf{n}|\,\in\,[k,\,k+1)} |V_\mathbf{n}|^2}{\Delta k_\mathrm{int}}
\end{equation}
Physical wavenumbers are $k = k_\mathrm{int}\cdot 2\pi/L$ in units of
$\mathrm{AU}^{-1}$, and the power spectral density is normalized by the voxel
volume $(\Delta x)^3$ so that $E(k)$ carries units of
$(\mathrm{km\,s^{-1}})^2\,\mathrm{AU}^3$.

The injection scale is marked at $k_\mathrm{inj} = 2\pi/r_\mathrm{disk}$ and the
dissipation scale at $k_\mathrm{diss} = 2\pi/\langle h\rangle$, where
$r_\mathrm{disk}$ is the 90th-percentile cylindrical radius of disk-classified
particles and $\langle h\rangle$ is their mean SPH smoothing length. A power-law
fit is performed in the inertial range $k_\mathrm{inj} \leq k \leq k_\mathrm{diss}$.
Reference slopes $k^{-5/3}$ (Kolmogorov incompressible turbulence),
$k^{-2}$ (Burgers highly compressible turbulence), and $k^{-3}$ (Kraichnan
2D enstrophy cascade) are overlaid for comparison.

% ─────────────────────────────────────────────────────────────────────────────

\section{Magnetic Field Analysis}

\subsection*{Magnetic field decomposition}

The magnetic field vector $\mathbf{B}$ of each gas particle is rotated into the
disk frame (aligned with the angular momentum axis $\hat{L}$) and decomposed into
cylindrical components:
\begin{equation}
    B_r = \mathbf{B} \cdot \hat{e}_r, \qquad
    B_\phi = \mathbf{B} \cdot \hat{e}_\phi, \qquad
    B_z = \mathbf{B} \cdot \hat{L}
\end{equation}
where $\hat{e}_r$ and $\hat{e}_\phi$ are the radial and azimuthal unit vectors in
the disk plane. The poloidal component is defined as
$B_\mathrm{pol} = \sqrt{B_r^2 + B_z^2}$, the toroidal component is $B_\phi$,
and the total magnitude is $|\mathbf{B}| = \sqrt{B_r^2 + B_\phi^2 + B_z^2}$.

\subsection*{Projected maps}

Face-on and edge-on maps of each magnetic field component are produced as
mass-weighted Meshoid projections:
\begin{equation}
    \langle B_\alpha \rangle(x,y)
    = \frac{\sum_i m_i\,B_{\alpha,i}\,W}{\sum_i m_i\,W}
\end{equation}
This is a signed projection (preserving field direction) for $B_z$, $B_r$, and
$B_\phi$, and unsigned for $|\mathbf{B}|$.

\subsection*{Magnetic energy}

The magnetic energy of the disk is computed as a particle-based volume integral:
\begin{equation}
    E_\mathrm{mag} = \sum_{i \in \mathrm{disk}} \frac{|\mathbf{B}_i|^2}{8\pi}
    \frac{m_i}{\rho_i}
\end{equation}
where $m_i/\rho_i$ approximates the volume element of each SPH particle. The
sum runs over all particles identified as disk members (Section~1). The magnetic
field is in Gaussian units (as stored by GIZMO after unit conversion), giving
$E_\mathrm{mag}$ in erg.

\subsection*{Mass-to-flux ratio}

The dimensionless mass-to-flux ratio $\mu_\Phi$ measures the ratio of the disk's
gravitational binding to its magnetic support, normalized to the critical value
for magnetic support against collapse \citep{mouschovias1976,nakano1978}:
\begin{equation}
    \mu_{\Phi,b} = \frac{2\pi\sqrt{G}\,\Sigma_b}{|B_{z,b}|}
\end{equation}
where $\Sigma_b$ is the surface density and $|B_{z,b}|$ is the mass-weighted mean
of the absolute vertical field component in annulus $b$. A value $\mu_\Phi > 1$
indicates that the disk is magnetically supercritical (gravity dominates over
magnetic support).

A pixel-by-pixel face-on map is computed analogously:
\begin{equation}
    \mu_\Phi(x,y) = \frac{2\pi\sqrt{G}\,\Sigma(x,y)}{|B_z(x,y)|}
\end{equation}
using the Meshoid-projected $\Sigma$ and $B_z$ maps. Both face-on and edge-on
projections are produced for comparison.

% ─────────────────────────────────────────────────────────────────────────────

\section*{Caveats}

\begin{itemize}
    \item \textbf{Mean molecular weight.} The temperature formula uses $\mu = 1$
          (neutral atomic hydrogen). At the densities found in the disk,
          primordial gas may be partially or fully molecular, in which case
          $\mu \approx 1.22$ (fully atomic) or $\mu \approx 2.46$ (fully
          molecular) would be more appropriate.

    \item \textbf{Adiabatic index.} The sound speed is derived assuming
          $\gamma = 5/3$ throughout. Molecular hydrogen has $\gamma_\mathrm{eff}
          \approx 7/5$ and radiation cooling further modifies the effective EOS.
          Depending on how GIZMO stores the effective equation of state for this
          run, $c_s$ derived from internal energy may already reflect the true
          EOS, making this a non-issue.
\end{itemize}

\begin{thebibliography}{9}
\bibitem{grudic2021}
    Grudić, M.~Y. et al., 2021, \textit{MNRAS}, 506, 2199
    (\texttt{starforge\_tools})
\bibitem{mouschovias1976}
    Mouschovias, T.~Ch. \& Spitzer, L., 1976, \textit{ApJ}, 210, 326
\bibitem{nakano1978}
    Nakano, T. \& Nakamura, T., 1978, \textit{PASJ}, 30, 671
\end{thebibliography}

\end{document}
